\chapter{Conclusiones}
\label{cap:conclusiones}

\section{Respuesta a las preguntas de investigación}

La evaluación de diversos modelos para predecir el volumen de capturas pesqueras en la temporada arrojó los siguientes resultados:

ARIMA, utilizado inicialmente en el contexto de todas las especies antes de ser descartado para la langosta, mostró un desempeño desfavorable y quedó fuera de la selección final.

Entre los modelos restantes, LSTM destacó como el más prometedor, exhibiendo el menor error absoluto medio y un error en la suma total de capturas de la temporada menor a los demás. Sin embargo, al visualizar los valores predichos frente a los reales, se observó que el ajuste de LSTM no seguía la tendencia de los datos reales. En contraste, XGBoost y LGBM se aproximaron mejor a dicha tendencia, a pesar de tener errores más grandes. Prophet y ARIMA, con ajustes simples que no siguen la tendencia ni se acercan a los valores reales, resultaron ininterpretables. 

El desempeño de los dos mejores modelos (LSTM y XGBoost) combinado superó a todos los demás, con un error porcentual en la suma total de pesca por temporada de 12.9\%. 

La escasez de datos de entrenamiento juega un papel importante en el mal desempeño de todos los modelos propuestos, por lo que no podemos pensar en modelos más complejos para mejorar los resultados, sino en aumentar la cantidad de datos y utilizar el mismo método para mejorar los resultados actuales.

En ese sentido, alcanzar un modelo aplicable a varias zonas de pesca y a varias especies depende de la disponibilidad de los datos para esos contextos; difícilmente los resultados de este trabajo nos acercan al resultado deseado, pero nos arrojan pistas de qué debemos hacer para alcanzarlo.

El trabajo desarrollado en esta tesis refuerza y precisa esas conclusiones iniciales. En primer lugar, la caída atípica de la temporada 2021-2022 —que el planteamiento inicial no lograba explicar y que sesgaba la evaluación de los modelos— se atribuye ahora a las olas de calor marinas de 2019-2021, un mecanismo ambiental documentado para Baja California \citep{villasenor2024} y reproducible con el índice de MHW que se incorporó al pipeline. De aquí se desprende una recomendación concreta: incluir siempre una variable de anomalía térmica de largo plazo (idealmente un índice de MHW) entre las covariables, pues es la señal que permite anticipar los años atípicos.

En segundo lugar, los experimentos con modelos multivariados y con un modelo global multi-especie confirman que el factor limitante sigue siendo la cantidad de temporadas y de series disponibles, no la sofisticación del modelo. La evidencia de transferencia positiva dentro de una misma especie entre unidades económicas es alentadora, y el enfoque jerárquico de \emph{partial pooling} que señalaba el camino a seguir quedó validado: al normalizar la variable objetivo por serie con $\log(1+y)$ —evitando que las series de mayor volumen dominen el ajuste— el modelo global pasa a ganar o empatar contra los modelos específicos en 4 de 5 series, rescatando en particular a las especies de escala pequeña. Escalar a más unidades económicas y acumular más temporadas —junto con el índice de MHW como covariable estándar y el objetivo en escala logarítmica— constituye la ruta más prometedora hacia un modelo operativo y transferible.

En tercer lugar, y como respuesta directa a esa escasez de datos, el pronóstico se reformuló en términos probabilísticos: una regresión cuantílica conformalizada sobre el modelo global entrega intervalos con cobertura cercana a la nominal y estable incluso durante las olas de calor marinas. Para una cooperativa pesquera, un rango calibrado de captura esperada es más accionable ---y más honesto--- que un único número puntual, y es el producto operativo concreto de este trabajo (Capítulo~\ref{cap:producto}).

Finalmente, la hipótesis que recorre todo el trabajo ---que el cuello de botella es el volumen de datos y no la sofisticación del modelo--- dejó de ser una conjetura y se puso a prueba. Al sumar diecinueve cooperativas de langosta (El Rosario, la costa norte, la península de Vizcaíno, el Pacífico de BCS y Bahía Magdalena, de 2 a 21 series) y desbloquear las temporadas 2022-2026 mediante la unión con los avisos de cosecha de CONAPESCA, la langosta de San Quintín ---cuya calibración ni la selección de variables por SHAP ni la optimización bayesiana de los cuantílicos habían hecho fiable--- pasó a estar bien calibrada de forma \emph{estable} (cobertura $\sim$96\%) al llevar el bache post-ola-de-calor de 2021-2022 dentro del entrenamiento, e \emph{invariante} a añadir zonas al pool; con el bache fuera del entrenamiento su cobertura era irreproducible, oscilando entre 40.6\% y 99.9\% según la composición del pool. La expansión también fijó el límite de esa transferencia: sumar zonas vecinas y de régimen térmico parecido afina el pronóstico, pero el extremo cálido del rango (Bahía Magdalena, $\sim$24$^\circ$N) ya no lo hace ---la palanca de fondo sigue siendo acumular temporadas, no zonas. Es la demostración más concreta de la tesis de fondo: la ruta hacia un modelo operativo y transferible pasa por acumular temporadas y unidades económicas ---junto con el índice de MHW como covariable estándar y el objetivo en escala logarítmica--- y no por más maquinaria estadística.

Esa hipótesis obligó también a revisar la conclusión del propio planteamiento inicial. Al re-entrenar sobre el conjunto unido los modelos que allí ganaban ---la LSTM y el ensamble XGBoost+LSTM--- su ventaja se disuelve: con tres temporadas de entrenamiento encabezan la tabla, pero por amortiguar la serie (razón de desviaciones 1.58 frente a 2.37 del XGBoost, con la misma correlación), no por seguir su forma; y con siete temporadas quedan cerca de cinco veces por detrás del XGBoost podado por SHAP, con errores de suma de temporada de $+626\%$ y $+704\%$. El error del 12.9\% que la primera iteración reportaba para el ensamble era, entonces, un artefacto de aquella partición corta y no una propiedad del modelo.

Como prueba de techo de esa afirmación se enfrentó al producto contra un Transformer de fusión temporal (TFT) con idénticas covariables y envuelto en la misma corrección conformal. El TFT resultó marginalmente algo más nítido (CRPS ligeramente a su favor en ambos cortes) pero \emph{mal calibrado}: sub-cubre en los dos cortes (70.8\% y 79.1\% para un intervalo nominal del 90\%, frente a la sobre-cobertura controlada del 96.7\% y 95.8\% del XGBoost), su diagnóstico por serie cambia de un extremo al otro al re-entrenar sobre otro pool, arrastra desbordes numéricos ocasionales y no gana nada al entrenar más épocas. Que la arquitectura de vanguardia no supere al gradient boosting con estos datos, mientras que sumar temporadas sí volvió fiable la calibración, es la confirmación más nítida de que la complejidad del modelo no es el factor limitante: lo es el volumen de datos.
% % Acknowledgements should go at the end, before appendices and references

% \acks{We would like to acknowledge support for this project
% from the National Science Foundation (NSF grant IIS-9988642)
% and the Multidisciplinary Research Program of the Department
% of Defense (MURI N00014-00-1-0637). }

% % Manual newpage inserted to improve layout of sample file - not
% % needed in general before appendices/bibliography.


\section{Limitaciones}
\label{sec:limitaciones}

Las limitaciones se enuncian en orden de gravedad para el uso del pronóstico.

\paragraph{La captura no es abundancia.} La variable objetivo es el volumen desembarcado, que
es el producto de la abundancia del recurso \emph{y} del esfuerzo pesquero: número de
embarcaciones, días efectivos de salida, precio de venta, decisiones de la cooperativa. No se
dispone de datos de esfuerzo, de modo que no es posible normalizar a captura por unidad de
esfuerzo (CPUE), que es la medida estándar en evaluación pesquera. Una parte de lo que el
modelo aprende es, por tanto, comportamiento humano y no dinámica poblacional, y una caída
puede reflejar menos salidas y no menos langosta. Esta es la limitación más importante para
interpretar los resultados en clave biológica.

\paragraph{Pocas temporadas de calibración.} Es la limitación que el propio trabajo documenta
como el cuello de botella. Con una a dos temporadas de calibración por serie, un cambio de
régimen queda fuera de la distribución aprendida: con el corte de 2020 la cobertura de la serie
insignia resulta irreproducible, oscilando según la composición del pool. El corte de 2024 es
el fiable, y lo es porque el bache está dentro del entrenamiento.

\paragraph{Zonas aproximadas por rectángulos.} La oceanografía se agrega sobre un rectángulo
costero y no sobre el polígono TURF real de cada cooperativa, que no es de acceso público. A la
resolución de los productos satelitales la aproximación es razonable, pero las cooperativas que
comparten rectángulo comparten también, por construcción, sus covariables oceanográficas.

\paragraph{Empalme de fuentes sin periodo de solape.} La unión COBI--CONAPESCA se apoya en la
continuidad del traspaso de diciembre de 2021 a enero de 2022. No existe un periodo en el que
ambas fuentes reporten simultáneamente, así que no fue posible validar exhaustivamente que las
escalas coincidan.

\paragraph{No se pronostica la oceanografía.} El horizonte útil llega hasta la última fecha con
covariables disponibles. El sistema es una herramienta de pronóstico calibrado sobre el periodo
con datos, no un oráculo de temporadas futuras arbitrarias; extenderlo exigiría alimentar
pronósticos oceanográficos, con su propia incertidumbre.

\paragraph{Dependencia de terceros.} Salvo los datos de COBI, obtenidos por convenio, las
fuentes son públicas pero externas: el trabajo depende de que Copernicus y CONAPESCA sostengan
sus procesos de publicación. Durante este trabajo, de hecho, el cliente \texttt{motuclient} de
Copernicus fue descontinuado y hubo que reescribir la ingesta.

\paragraph{Una sola especie con evidencia sólida.} Las conclusiones sobre transferencia entre
zonas están respaldadas para la langosta roja, con 21 series. Las de abulón y erizo tienen menos
series y menos temporadas, y sus resultados deben leerse como indicativos.

\section{Trabajo futuro}
\label{sec:trabajo_futuro}

\begin{enumerate}
  \item \textbf{Incorporar esfuerzo pesquero.} Es la mejora de mayor impacto científico:
  permitiría modelar CPUE y separar la señal ambiental del comportamiento de la flota. Requiere
  datos de salidas por embarcación, que las cooperativas sí registran.

  \item \textbf{Acumular temporadas.} Es la conclusión que atraviesa todo el trabajo. Cada
  temporada adicional vale más que cualquier refinamiento de modelo, y el sistema ya está
  construido para incorporarlas con un solo comando.

  \item \textbf{Conformal adaptativo por régimen.} Calibrar por separado los días con y sin ola
  de calor daría intervalos que se ensanchan automáticamente en régimen anómalo. Con una sola
  temporada de calibración por serie hoy no es robusto, pero lo será al acumular temporadas.

  \item \textbf{Polígonos TURF reales} y, con ellos, covariables oceanográficas distintas por
  cooperativa dentro de un mismo clúster.

  \item \textbf{Extender el catálogo de especies} a las que hoy tienen series cortas,
  aprovechando que el modelo global ya está diseñado para agruparlas.

  \item \textbf{Validar el producto con las cooperativas.} La utilidad real del intervalo
  calibrado solo puede evaluarse con quienes toman la decisión de compra de insumos; el
  siguiente paso natural es una temporada de uso acompañado.
\end{enumerate}
